\chapter{Kết luận}
\label{chap:conclusion}

\section{Tóm tắt}

Báo cáo đã trình bày trọn một vòng đời: từ kiến trúc năm tầng
(Chương~\ref{chap:implementation}) qua mô hình dữ liệu đa chỉ số đến phương
pháp đánh giá theo năm câu hỏi nghiên cứu. Toàn bộ hệ thống --- lớp tiếp nhận
MQTT, bộ đệm Kafka, engine Spark Structured Streaming, kho Parquet/
PostgreSQL, tầng phục vụ FastAPI/dashboard và tác vụ batch theo giờ được
kích hoạt bằng script --- chạy
trong Docker Compose trên một máy Linux với ngân sách bộ nhớ live mode
khoảng 4.6~GiB.

Khung khái niệm veracity-aware event-time ở Chương~\ref{chap:background} ---
từ phân loại thời gian, lateness, watermark đến veracity vector và các bất
biến kế toán bản ghi --- đã trở thành hệ quy chiếu dùng xuyên suốt khi thiết
kế cơ chế validation và khi đọc kết quả đo. Mỗi bản tin là một vector
data-fusion gồm sáu chỉ số môi trường (nhiệt độ, độ ẩm, CO$_2$, áp suất, bụi
mịn PM2.5, ánh sáng); kiểm tra chất lượng và phát hiện cảnh báo đều được
tổng quát hoá theo bảng cấu hình per-metric thay vì viết cứng cho một đại
lượng; cảnh báo được tính ngay trong Q3, chia sẻ transaction với việc ghi
dữ liệu phục vụ.

Năm câu hỏi nghiên cứu đã đặt ra trong Chương~\ref{chap:introduction} được
ánh xạ vào phương pháp benchmark gồm bảy kịch bản và mười ba nhóm metric
thu thập từ query progress JSONL, Kafka và thống kê container. Một phép đo
runtime dài đã hoàn thành với run \texttt{load15-clean-20261002}: năm nguồn
phát trong bốn giờ, 50.760.004 sự kiện và 16.226.141.856 byte payload được
ghi nhận tại Bronze. Kết quả này xác nhận mốc 15 GB/4 giờ của cấu hình một
máy, cùng khả năng drain Kafka và ổn định Q1--Q3. Nó không xác nhận mốc
15 GB/3 giờ; mốc đó không nằm trong kế hoạch hiện tại. Các kết quả còn lại
được phân loại theo evidence: một số là kiểm thử fixture/định tính đã có,
phần đường cong saturation và mục tiêu 30 GB/4 giờ vẫn là công việc tiếp
theo:

\begin{itemize}
  \item \textbf{RQ-A:} phép tải dài đã đo được 3.523,76 event/s tổng,
        16,226 GB payload trong bốn giờ và P95 micro-batch Q1 là 422 ms.
        Đây là một điểm đo ổn định của pipeline; đường cong khi vượt capacity
        vẫn cần các workload rate-sweep riêng;
  \item \textbf{RQ-B:} cấu hình một broker, ba partition và
        \texttt{local[2]} đã xử lý được workload 15 GB/4 giờ, nhưng phép đo
        này chưa cô lập ảnh hưởng của partition count, số core hoặc số Bridge.
        Các kết luận scale-up trong phần tham chiếu vì vậy được giữ ở mức
        giả thuyết thiết kế;
  \item \textbf{RQ-C:} checkpoint trên bind mount và Kafka retention tạo seam
        phục hồi trong kiến trúc, còn kết luận recovery time cần phép dừng/
        khởi động riêng. Không suy ra recovery từ run dài sạch;
  \item \textbf{RQ-D:} fixture trước đó đã kiểm tra validation per-metric,
        quarantine, deduplication và late flag; run dài mới tắt fault
        injection nên chỉ xác nhận đường đi dữ liệu sạch, không xác nhận tỷ
        lệ lỗi hoặc chất lượng dữ liệu hỏng;
  \item \textbf{RQ-E:} Bronze và Silver đều ghi nhận 50.760.004 dòng trong
        run dài, nhưng Gold riêng cho prefix mới chưa chạy. Vì vậy tính xác
        định của batch rerun được bảo vệ bởi fixture Gold trước đó, chưa phải
        kết quả Gold của workload 15 GB/4 giờ.
\end{itemize}

Về phương pháp, báo cáo chủ trương demo nhỏ nhưng giữ nguyên các seam kiến
trúc: mọi ranh giới cần cho mở rộng ngang sau này (số broker, số partition,
số bản sao bridge, số nhân Spark) đều đã hiện diện trong cấu hình, nên việc
mở rộng về sau chỉ là thay đổi tham số chứ không phải viết lại pipeline.

\section{Hướng phát triển}

Các hướng phát triển được chia thành hai nhóm: những gì demo được trên một
máy (đã đưa vào kịch bản benchmark đối chứng ở Chương~\ref{chap:results}) và
những gì cần nhiều máy hơn hoặc chưa triển khai trong bản này.

\subsection{Mục tiêu tải tiếp theo: 30 GB trong 4 giờ}

Phép tải 15 GB/4 giờ đã cung cấp mốc thực nghiệm cho cấu hình một máy.
Mục tiêu tiếp theo là đưa ít nhất 30 GB payload tới Bronze trong cùng thời
lượng bốn giờ. Với kích thước payload trung bình 319,66 byte của run trước,
điều này tương đương khoảng 2,083 MB/s, 6.517 event/s tổng hoặc 1.303
event/s mỗi nguồn nếu vẫn giữ năm nguồn. Năng lực payload cần cao hơn
khoảng 1,85 lần so với phép đo hiện tại. Phép thử này chưa được chạy; trước
khi thực hiện cần xử lý ngân sách đĩa, giữ retention/evidence phù hợp và
đối soát riêng nguồn, Bridge, Kafka, Bronze, Silver và PostgreSQL.

\subsection{Mở rộng ingestion và xử lý song song (demo được trên 1 máy)}

Các kịch bản này đã được xác định trong phương pháp benchmark và có thể chạy
đối chứng trên máy mục tiêu mà không thay đổi kiến trúc cốt lõi:

\begin{itemize}
  \item Nhân bản N Bridge song song với \texttt{BRIDGE\_INSTANCE\_ID} duy nhất,
        đo throughput cộng dồn và xung đột idempotent producer khi thiếu
        \texttt{client.id} riêng;
  \item So sánh 1 topic gom chung vs N topic tách theo zone, đo overhead
        metadata broker;
  \item Suy rộng số partition (1 $\rightarrow$ 3 $\rightarrow$ 6), đo đường cong
        throughput Spark Q1 để tìm điểm bão hoà của \texttt{local[N]};
  \item Tăng số core xử lý (\texttt{local[2]} $\rightarrow$ \texttt{local[4]}),
        đo ảnh hưởng lên throughput tổng hợp của 4 query;
  \item Bật persistence Mosquitto, kill Bridge giữa chừng, đo khả năng giữ
        và redeliver message của broker.
\end{itemize}

\subsection{Kiểm chứng extension veracity thích nghi}

Triển khai chính sách FAST/CAUTIOUS tại seam đã nêu ở
Chương~\ref{chap:implementation} (controller đọc metric chất lượng per-metric
từ quality meta-stream, chọn lại tham số watermark/validation và khởi động
lại query giữa các epoch) cùng vòng đời công bố đầy đủ với ba trạng thái
PROVISIONAL, FINAL, CORRECTED và bước reconciliation batch. Sau đó chạy bốn
phép
đo ứng với H1--H4: (i) bộ ba retention--state--publication delay theo
$\theta$; (ii) kế toán chất lượng so với profile inject; (iii) sai lệch
aggregate trước/sau reconciliation; (iv) ablation fixed vs adaptive, có/không
hysteresis. Đây là phần biến các giả thuyết hiện tại thành kết quả đo được;
mọi tuyên bố về lợi ích adaptive đều phải chờ các phép đo này.

\subsection{Kubernetes và production deployment}

Khi pipeline đã chạy ổn định trên Docker Compose, bước tiếp theo là chuyển
sang Kubernetes. Toàn bộ stack đều có tài liệu chính thức cho lộ trình này:
Spark on Kubernetes, Strimzi cho Kafka; lúc đó việc điều phối batch bằng
script/cron được thay bằng Airflow (Helm chart) khi số DAG và phụ thuộc
giữa các tác vụ tăng lên. So với Docker
Compose, Kubernetes cung cấp rolling update, tự phục hồi cho Spark
driver/executor pod, và khả năng mở rộng theo metric (HPA).

\subsection{High availability}

Hai điểm cần xử lý để đạt HA ở mức production: (i) Kafka replication factor
$\geq 3$ với min in-sync replica $= 2$ (bản demo dùng RF=1, 1 broker KRaft);
(ii) PostgreSQL read replica hoặc Patroni cho HA tầng serving. Khi kết hợp với
Kubernetes, có thể thiết lập pod disruption budget để tránh mất nhiều replica
cùng lúc. Lưu ý rằng HA thật đòi hỏi tối thiểu hai máy vật lý --- demo trên
một máy chỉ chứng minh được seam kiến trúc, không chứng minh được khả năng
chịu lỗi phần cứng.

\subsection{Bảo mật}

Hệ thống mặc định vận hành trong network nội bộ để dễ debug (Mosquitto
anonymous/plaintext, Kafka không SASL). Hướng phát triển gồm: (i) bật TLS
cho MQTT (port 8883) và Kafka giữa broker và client; (ii) SASL/SCRAM hoặc
mTLS cho authentication; (iii) ACL theo topic MQTT và Kafka ACL để giới hạn
producer chỉ WRITE \texttt{sensor\_raw}, Spark chỉ READ, admin operation tách
riêng~\cite{kafka_security_overview}. Các seam bảo mật đã sẵn sàng ở mức cấu
hình (\texttt{mosquitto.conf}, Kafka listener properties) --- bật lên chỉ cần
thêm directive, không sửa logic Bridge hay Spark.

\subsection{Nâng cấp storage layer}

Bản demo dùng plain Parquet trên bind mount cục bộ. Hướng phát triển gồm:
(i) chuyển sang Delta Lake hoặc Apache Iceberg trên object storage (S3/MinIO)
để có ACID transaction, time travel, schema evolution và compaction tự động
(\texttt{OPTIMIZE}/\texttt{VACUUM}); (ii) sử dụng HDFS nếu cần replication
ở tầng tệp. Việc nâng cấp này chỉ thay đổi format writer/reader trong Spark,
không thay đổi logic Q1--Q3 hay mô hình Bronze/Silver/Gold.

\subsection{Quan sát vận hành tập trung}

Thay thế script \texttt{resource-report.sh} bằng cụm Prometheus (pull model)
+ JMX Exporter cho JVM của Kafka/Spark + Grafana dashboard/alerting. Các
metric cần thu thập đã được xác định rõ trong implementation hiện tại (query
progress JSONL, docker stats, Kafka end-offsets) --- chỉ cần thêm exporter,
không thay đổi pipeline.

\subsection{Edge / fog processing}

Khi yêu cầu latency cỡ mili-giây hoặc khi băng thông uplink hạn chế, một phần
xử lý có thể đẩy xuống edge gateway~\cite{iot_big_data_survey_2022,
iot_stream_latency_2021}. Bước đầu có thể chạy Spark Structured Streaming trên
edge node với micro-batch interval ngắn hơn, sau đó gửi aggregate về cloud.
Hướng này đòi hỏi thay đổi đáng kể về kiến trúc và nằm ngoài phạm vi đồ án
hiện tại.

\subsection{Mở rộng RQ và dataset}

Báo cáo sử dụng dataset tự sinh từ simulator đa nguồn. Hướng phát triển gồm:
(i) replay thêm các dataset IoT công khai khác (ví dụ Intel Berkeley
~\cite{intel_berkeley_dataset}) để kiểm tra tính tổng quát của cấu hình;
(ii) bổ sung RQ về \emph{energy efficiency} và \emph{cost per million events};
(iii) so sánh với một hệ thống tham chiếu khác (ví dụ Flink) để đánh giá tính
tương đối; (iv) tích hợp framework giám sát chất lượng stream-first như
Stream DaQ~\cite{stream_daq_2024} để mở rộng từ tuple-at-a-time check sang
window-context và dynamic-context check.

\section{Lời kết}

Đề tài này được chọn với quan điểm: không nên biến bài toán cảm biến IoT
thành một bài phân tích dữ liệu đơn lẻ, mà phải nhìn nhận nó như một bài toán
\emph{hệ thống} --- nơi các đặc trưng Volume, Velocity, Veracity và tính phân
tán mới là đối tượng nghiên cứu chính. Bằng cách chọn stack mã nguồn mở phổ
biến, có tài liệu chính thức và có cộng đồng, cùng một khung khái niệm
veracity-aware event-time có thể kiểm chứng và một phương pháp benchmark có
cơ sở học thuật~\cite{kafka_patterns_benchmark_2025, stream_daq_2024}, báo cáo
cung cấp một mốc đo thực tế 15 GB/4 giờ cho nền tảng một máy, đồng thời
giữ mục tiêu 30 GB/4 giờ như bước đánh giá kế tiếp trên cùng pipeline.

</content>